\begin{obereflection}
\textbf{Latihan 2.1}
\begin{enumerate}
    \item Mengapa \textit{chosen-plaintext attack} dianggap lebih berbahaya
          daripada \textit{ciphertext-only attack}? Berikan satu contoh skenario.
    \item Berikan contoh \textit{replay attack} pada sistem perbankan daring dan
          jelaskan bagaimana \textit{nonce} atau \textit{timestamp} mencegahnya.
    \item Susun kelima jenis serangan berikut dari yang paling lemah ke yang
          paling kuat bagi penyerang, lalu jelaskan alasannya:
          \textit{chosen-ciphertext}, \textit{known-plaintext},
          \textit{ciphertext-only}, \textit{adaptive chosen-plaintext},
          \textit{chosen-plaintext}.
    \item Seorang penyerang mengetahui bahwa komunikasi antara dua kantor hukum
          meningkat menjelang suatu sidang, meskipun isi pesannya tersandi.
          Serangan jenis apa yang ia lakukan?
    \item Sebuah kartu pintar menyimpan kunci rahasia dan tidak dapat dibongkar
          secara fisik. Sebutkan dua serangan saluran samping yang masih dapat
          dipakai penyerang.
    \item Sebutkan tiga syarat sebuah algoritma disebut aman secara komputasi.
\end{enumerate}
\end{obereflection}
